\documentclass[10pt,a4paper]{amsart}
\usepackage{amsmath,amssymb,amsthm,graphicx}
% Retained minimal preamble for the Stacks Project chapter flat.tex (More on Flatness).
% Verbatim excerpt of the chapter preamble: Xy-pic setup (native preamble lines 14-18) and the native
% theorem declarations with the mathematical macros (native preamble lines 154-217). Copied mechanically;
% nothing added, removed or redefined. The wrapper class and the amsmath/amssymb/amsthm/graphicx base
% packages are supplied by the assembler, so the class/externaldocument part of the native preamble is omitted.
% For commutative diagrams we use Xy-pic
\usepackage[all]{xy}

% We use 2cell for 2-commutative diagrams.
\xyoption{2cell}
% Theorem environments.
%
\theoremstyle{plain}
\newtheorem{theorem}[subsection]{Theorem}
\newtheorem{proposition}[subsection]{Proposition}
\newtheorem{lemma}[subsection]{Lemma}

\theoremstyle{definition}
\newtheorem{definition}[subsection]{Definition}
\newtheorem{example}[subsection]{Example}
\newtheorem{exercise}[subsection]{Exercise}
\newtheorem{situation}[subsection]{Situation}

\theoremstyle{remark}
\newtheorem{remark}[subsection]{Remark}
\newtheorem{remarks}[subsection]{Remarks}

\numberwithin{equation}{subsection}

% Macros
%
\def\lim{\mathop{\mathrm{lim}}\nolimits}
\def\colim{\mathop{\mathrm{colim}}\nolimits}
\def\Spec{\mathop{\mathrm{Spec}}}
\def\Hom{\mathop{\mathrm{Hom}}\nolimits}
\def\Ext{\mathop{\mathrm{Ext}}\nolimits}
\def\SheafHom{\mathop{\mathcal{H}\!\mathit{om}}\nolimits}
\def\SheafExt{\mathop{\mathcal{E}\!\mathit{xt}}\nolimits}
\def\Sch{\mathit{Sch}}
\def\Mor{\mathop{\mathrm{Mor}}\nolimits}
\def\Ob{\mathop{\mathrm{Ob}}\nolimits}
\def\Sh{\mathop{\mathit{Sh}}\nolimits}
\def\NL{\mathop{N\!L}\nolimits}
\def\CH{\mathop{\mathrm{CH}}\nolimits}
\def\proetale{{pro\text{-}\acute{e}tale}}
\def\etale{{\acute{e}tale}}
\def\QCoh{\mathit{QCoh}}
\def\Ker{\mathop{\mathrm{Ker}}}
\def\Im{\mathop{\mathrm{Im}}}
\def\Coker{\mathop{\mathrm{Coker}}}
\def\Coim{\mathop{\mathrm{Coim}}}

% Boxtimes
%
\DeclareMathSymbol{\boxtimes}{\mathbin}{AMSa}{"02}

%
% Macros for moduli stacks/spaces
%
\def\QCohstack{\mathcal{QC}\!\mathit{oh}}
\def\Cohstack{\mathcal{C}\!\mathit{oh}}
\def\Spacesstack{\mathcal{S}\!\mathit{paces}}
\def\Quotfunctor{\mathrm{Quot}}
\def\Hilbfunctor{\mathrm{Hilb}}
\def\Curvesstack{\mathcal{C}\!\mathit{urves}}
\def\Polarizedstack{\mathcal{P}\!\mathit{olarized}}
\def\Complexesstack{\mathcal{C}\!\mathit{omplexes}}
% \Pic is the operator that assigns to X its picard group, usage \Pic(X)
% \Picardstack_{X/B} denotes the Picard stack of X over B
% \Picardfunctor_{X/B} denotes the Picard functor of X over B
\def\Pic{\mathop{\mathrm{Pic}}\nolimits}
\def\Picardstack{\mathcal{P}\!\mathit{ic}}
\def\Picardfunctor{\mathrm{Pic}}
\def\Deformationcategory{\mathcal{D}\!\mathit{ef}}

% --- appended 2026-10-12 (Chief): the project own definitions of the reference/slogan/history
% display environments, copied mechanically from the shared Stacks preamble lines 8-12
% (\usepackage{verbatim} plus the three \newenvironment lines, which the chapter preamble keeps
% verbatim from the project preamble). Nothing added, removed or redefined: these environments
% carry no mathematical content and the project itself typesets them as comments. Without these
% three lines a counted statement whose own text uses them fails with Environment slogan undefined.
\usepackage{verbatim}
\newenvironment{reference}{\comment}{\endcomment}
\newenvironment{slogan}{\comment}{\endcomment}
\newenvironment{history}{\comment}{\endcomment}
\usepackage{cleveref}
\crefname{originallemma}{original lemma}{original lemmas}
\Crefname{originallemma}{Original lemma}{Original lemmas}
\makeatletter
\newlabel{algebra-lemma-colimit-universally-injective-ML}{{0AS7}{source}{}{sourcepointer-1}{}}
\newlabel{algebra-lemma-colimit-universally-injective-ML@cref}{{[originallemma][0][]0AS7}{[1][0][]source}{}{}{}}
\makeatother
\crefname{originallemma}{original lemma}{original lemmas}
\Crefname{originallemma}{Original lemma}{Original lemmas}
\makeatletter
\newlabel{lemma-flat-finite-type-affine-neighbourhood-projective}{{05MF}{source}{}{sourcepointer-2}{}}
\newlabel{lemma-flat-finite-type-affine-neighbourhood-projective@cref}{{[originallemma][0][]05MF}{[1][0][]source}{}{}{}}
\makeatother
\crefname{originallemma}{original lemma}{original lemmas}
\Crefname{originallemma}{Original lemma}{Original lemmas}
\makeatletter
\newlabel{lemma-universally-injective-local}{{05FQ}{source}{}{sourcepointer-3}{}}
\newlabel{lemma-universally-injective-local@cref}{{[originallemma][0][]05FQ}{[1][0][]source}{}{}{}}
\makeatother
\title{A flat finite module over a henselian local base is a filtered colimit of free modules with injective reductions}
\author{The Stacks Project Authors}
\date{stacks-project:flat.tex:local-snapshot-20261011; GNU FDL 1.2 or later}
\begin{document}
\maketitle
\noindent\textit{Editorial note.} All delivered mathematical text is the author's own, copied from the pinned Stacks Project chapter file flat.tex; no statement or proof step was rewritten, reordered or removed.\par\medskip
\noindent\textit{Editorial note.} Preamble provenance: the chapter's own retained minimal preamble (sources/stacks-flat/retained-preamble.tex, lines 1-85) is a verbatim excerpt of the native Xy-pic setup and of the native theorem declarations and mathematical macros, followed by the project's own reference/slogan/history display-environment definitions copied mechanically from the shared preamble; the wrapper class and the amsmath/amssymb/amsthm/graphicx base packages are supplied by the assembler, and extra\_packages is left empty because everything needed is already in the retained preamble. No macro was invented or redefined.\par\medskip
\noindent\textit{Editorial note.} Same-chapter pointers are located at their \textbackslash{}label lines inside the pinned chapter file; cross-chapter pointers are located in the project's own tag list staged in the snapshot (sources/stacks-chow/tags.txt), at the verbatim lines carrying their canonical tags.\par\medskip
\noindent\textit{Allowed standard prerequisites (editorial).}\begin{itemize}
\item Projective neighbourhoods of a finite module flat at a point, in the henselian case. Same chapter flat.tex, label lemma-flat-finite-type-affine-neighbourhood-projective, invoked at proof line 5779. Gives a fundamental system of affine open neighbourhoods of the point on which the global sections of the module are free over the base, which produces the filtered system of free modules.
\item Universal injectivity of the transition maps, and the Mittag-Leffler criterion for a colimit of universally injective maps. Same chapter flat.tex, label lemma-universally-injective-local (proof line 5791); Algebra, label algebra-lemma-colimit-universally-injective-ML (proof line 5794). The first shows that the structure maps of the system are universally injective over the base, and the second converts that into the Mittag-Leffler property of the colimit.
\end{itemize}\medskip
\ifdefined\hypertarget\hypertarget{sourcepointer-1}{}\fi
\noindent\textit{Source reference.} Original lemma 0AS7: sources/stacks-chow/tags.txt, source lines 13570--13570. Navigation pointer for \textbackslash{}ref\{algebra-lemma-colimit-universally-injective-ML\} used inside the counted span (line 5794); the argument that consumes it is retained in the span.\par\medskip
\ifdefined\hypertarget\hypertarget{sourcepointer-2}{}\fi
\noindent\textit{Source reference.} Original lemma 05MF: sources/stacks-flat/pinned-original-source.tex, source lines 5728--5728. Navigation pointer for \textbackslash{}ref\{lemma-flat-finite-type-affine-neighbourhood-projective\} used inside the counted span (line 5779); the argument that consumes it is retained in the span.\par\medskip
\ifdefined\hypertarget\hypertarget{sourcepointer-3}{}\fi
\noindent\textit{Source reference.} Original lemma 05FQ: sources/stacks-flat/pinned-original-source.tex, source lines 1946--1946. Navigation pointer for \textbackslash{}ref\{lemma-universally-injective-local\} used inside the counted span (line 5791); the argument that consumes it is retained in the span.\par\medskip
\begin{lemma}
\label{lemma-flat-finite-type-local-colimit-free}
Let $A \to B$ be a local ring map of local rings which is essentially of
finite type. Let $N$ be a finite $B$-module which is flat as an $A$-module.
If $A$ is henselian, then $N$ is a filtered colimit
$$
N = \colim_i F_i
$$
of free $A$-modules $F_i$ such that all transition maps
$u_i : F_i \to F_{i'}$ of the system induce injective maps
$\overline{u}_i : F_i/\mathfrak m_AF_i \to F_{i'}/\mathfrak m_AF_{i'}$.
Also, $N$ is a Mittag-Leffler $A$-module.
\end{lemma}

\begin{proof}
We can find a morphism of finite type $X \to S = \Spec(A)$
and a point $x \in X$ lying over the closed point $s$ of $S$ and a finite
type quasi-coherent $\mathcal{O}_X$-module $\mathcal{F}$ such that
$\mathcal{F}_x \cong N$ as an $A$-module. After shrinking $X$
we may assume that each point of $\text{Ass}_{X_s}(\mathcal{F}_s)$ specializes
to $x$. By
Lemma \ref{lemma-flat-finite-type-affine-neighbourhood-projective}
we see that there exists a fundamental system of affine open neighbourhoods
$U_i \subset X$ of $x$ such that $\Gamma(U_i, \mathcal{F})$ is
a free $A$-module $F_i$. Note that if $U_{i'} \subset U_i$, then
$$
F_i/\mathfrak m_AF_i = \Gamma(U_{i, s}, \mathcal{F}_s)
\longrightarrow
\Gamma(U_{i', s}, \mathcal{F}_s) = F_{i'}/\mathfrak m_AF_{i'}
$$
is injective because a section of the kernel would be supported at
a closed subset of $X_s$ not meeting $x$ which is a contradiction
to our choice of $X$ above. Since the maps $F_i \to F_{i'}$ are
$A$-universally injective (Lemma \ref{lemma-universally-injective-local})
it follows that $N$ is
Mittag-Leffler by
Algebra, Lemma \ref{algebra-lemma-colimit-universally-injective-ML}.
\end{proof}

\end{document}
